% !TEX program = lualatex
\documentclass[12pt,a4paper]{article}
\usepackage{graphicx}
\usepackage{fontspec}
\setmonofont{Latin Modern Mono}
\usepackage{xcolor}
\usepackage{listings}
\usepackage{geometry}
\usepackage{fancyhdr}
\usepackage{hyperref}
\usepackage{booktabs} 
\usepackage{float}    
\usepackage{amsmath}

% =======================================================
%  USER CONFIGURATION
% =======================================================
\newcommand{\labNo}{Lab 10}
\newcommand{\labTopic}{PL/pgSQL: Anonymous Blocks}
\newcommand{\courseCode}{CSE 4308}
\newcommand{\courseTitleShort}{DBMS Lab}
\newcommand{\courseTitleLong}{DATABASE MANAGEMENT SYSTEMS LAB}
\newcommand{\instructorName}{Samnun Azfar, Jr. Lecturer, CSE\\ Samia Nawsheen, Jr. Lecturer, CSE}
\newcommand{\academicYear}{AY25-26}

% --- SOLUTION TOGGLE ---
% Change \showsolutionstrue to \showsolutionsfalse to generate 
% the student version (hidden answers).
\newif\ifshowsolutions
% \showsolutionstrue 
\showsolutionsfalse
% =======================================================

\geometry{margin=1in}

% -
% Header and Footer Configuration
% -
\pagestyle{fancy}
\fancyhf{}
\lhead{\labNo}                        % Uses variable
\rhead{\courseCode: \courseTitleShort} % Uses variables

% Centered footer with page number
\cfoot{\thepage}
% Ensure header/footer width matches text block
\setlength{\headwidth}{\textwidth}

% -
% Code Listing Style
% -
\definecolor{codegreen}{rgb}{0,0.6,0}
\definecolor{codegray}{rgb}{0.5,0.5,0.5}
\definecolor{codepurple}{rgb}{0.58,0,0.82}
\definecolor{backcolour}{rgb}{0.95,0.95,0.92}

\lstdefinestyle{mystyle}{
    backgroundcolor=\color{backcolour},
    commentstyle=\color{codegreen},
    keywordstyle=\color{magenta},
    numberstyle=\tiny\color{codegray},
    stringstyle=\color{codepurple},
    basicstyle=\ttfamily\footnotesize,
    breaklines=true,
    captionpos=b,
    keepspaces=true,
    numbers=left,
    numbersep=5pt,
    tabsize=2,
    frame=single,
    numbers=none,
    columns=fullflexible,
    showspaces=false,
    showstringspaces=false,
    showtabs=false,
}

\lstset{style=mystyle}

\begin{document}

% --
% TITLE PAGE
% --
\begin{titlepage}
    \centering
    \vspace*{2\baselineskip}

    \rule{14cm}{2pt}\vspace*{-0.8\baselineskip}\vspace*{1.5pt}
    \rule{14cm}{0.5pt}\vspace{2\baselineskip}

    {\Huge\bfseries\scshape \labNo\par} % Uses variable
    \vspace{3mm}
    {\Large \labTopic} % Uses variable
    \vspace{3mm}

    \rule{14cm}{0.5pt}\vspace*{-0.7\baselineskip}\vspace{2.5pt}
    \rule{14cm}{2pt}

    \vspace{2\baselineskip}

    {\Large\bfseries \courseCode\par} % Uses variable
    \vspace{1\baselineskip}
    {\large \courseTitleLong\par} % Uses variable

    \vfill
    {\normalsize Prepared by:\\ \instructorName\par} % Uses variable
    {\normalsize Academic Year: \academicYear\par} % Uses variable
\end{titlepage}

\newpage
\tableofcontents
\newpage

% =======================================================
% SECTION: LAB WORKFLOW
% =======================================================
\section{Lab Workflow}



\textit{Repeat this subsection for every single lab.}

\begin{enumerate}
\item Connect to the internet by running nmtui in terminal, add a dsl connection, then go back and activate that dsl connection. Make sure that your net id is not being used in any other pc or router, if you face any issue, use any other net id from your friends which is not being used.

\item Download all the files posted in the google classroom.

\item Disconnect from the \textbf{dsl} connection via nmtui, connect to the ethernet/wired connection via nmtui. If no such connection exists, create an \textbf{ethernet} connection  and then activate.

    \item \textbf{Create Your Role and Database}

Inside the PostgreSQL shell, create your own role (user) and database.

1. In the Linux terminal:
\begin{lstlisting}[language=bash]
psql
\end{lstlisting}

Expected prompt:
\begin{lstlisting}
cse=#
\end{lstlisting}

2. Create Your Role and Database

Copy and run the following commands inside PostgreSQL:
\begin{lstlisting}[language=SQL]
-- 1. Terminate existing database connections
SELECT pg_terminate_backend(pid)
FROM pg_stat_activity
WHERE datname = 'cse4308_12345'
  AND pid <> pg_backend_pid();

-- 2. Remove existing database
DROP DATABASE IF EXISTS cse4308_12345;

-- 3. Remove existing role
DROP ROLE IF EXISTS student_12345;

-- 4. Create role
CREATE ROLE student_12345 LOGIN;

-- 5. Set password
ALTER ROLE student_12345
WITH PASSWORD 'student12345';

-- 6. Create database
CREATE DATABASE cse4308_12345
OWNER student_12345;

-- 7. Grant privileges
GRANT ALL PRIVILEGES
ON DATABASE cse4308_12345
TO student_12345;
\end{lstlisting}

3. Exit PostgreSQL:

   \begin{lstlisting}[language=SQL]
   \q
   \end{lstlisting}

\item \textbf{Connect Using Your Newly Created User}

1. Connect to Your Database:

   Use the following command to log in as the \texttt{student\_12345} user:

   \begin{lstlisting}[language=SQL]
   psql -h localhost -U student_12345 -d cse4308_12345
   \end{lstlisting}

2. Verify the Connection:

   Verify that you are connected to the correct database with the following command:

   \begin{lstlisting}[language=SQL]
   \conninfo
   \end{lstlisting}

   Expected output:

   \begin{lstlisting}
   Database: cse4308_12345, User: student_12345
   \end{lstlisting}

3. Exit PostgreSQL:

   \begin{lstlisting}[language=SQL]
   \q
   \end{lstlisting}

\item \textbf{Load the Lab SQL File (MUST be run outside psql)}

Load the SQL file that sets up the tables and inserts initial data.

1. Load the SQL File:

   Run the following command to load the lab SQL file:

   \begin{lstlisting}[language=SQL]
   psql -h localhost -U student_12345 -d cse4308_12345 -f /home/cse/Downloads/lab_10_setup.sql
   \end{lstlisting}

   This will:
   - Drop old tables (if they exist).
   - Create new tables.
   - Insert initial rows of data.

\item \textbf{Reconnect and Verify Tables}

After loading the SQL file, reconnect to PostgreSQL to verify that the tables have been created.

1. Reconnect to Your Database:

   \begin{lstlisting}[language=SQL]
   psql -h localhost -U student_12345 -d cse4308_12345
   \end{lstlisting}

2. Check Tables:

   Verify the tables have been created by listing them:

   \begin{lstlisting}[language=SQL]
   \dt
   \end{lstlisting}

3. View Table Data:

   Verify the data by running the following SQL queries:

   \begin{lstlisting}[language=SQL]
   SELECT * FROM <table_name>;
   \end{lstlisting}

\item Open a text editor and save it in the appropriate submission name format \\ \texttt{fullid\_CSE4308\_L10\_T1.sql}. You will use this text editor to write your solutions, make sure to save it after writing each solution. The structure of this file will be:
\begin{lstlisting}[language=SQL]
-- Task 1
INSERT INTO students (student_id, name, major, gpa)
VALUES (104, 'David', 'ME', 3.60);
--Task 2
SELECT name, gpa FROM students WHERE major = 'CSE';
\end{lstlisting}
\end{enumerate}
=======================================================
% SECTION: DATABASE SCHEMA
% =======================================================
\section{Database Schema \& Data}
\textbf{Instructions:} Save this code as \texttt{lab\_10\_setup.sql} and run it using \texttt{\textbackslash i lab\_10\_setup.sql} in your psql terminal. This script generates a dataset of 1,000 employees with varying salaries and departments to simulate a real-world scenario.

\begin{lstlisting}[language=SQL, caption={lab\_10\_setup.sql}]
-- 1. CLEANUP
DROP TABLE IF EXISTS employees;

-- 2. CREATE TABLE
CREATE TABLE employees (
    employee_id SERIAL PRIMARY KEY,
    first_name  VARCHAR(50),
    last_name   VARCHAR(50),
    salary      NUMERIC(10, 2),
    department  VARCHAR(50)
);

-- 3. INSERT KEY PERSONNEL (Deterministic Data)
-- We insert these specific people so we can always find them in queries
INSERT INTO employees (first_name, last_name, salary, department) VALUES 
('Alice', 'Director', 25000.00, 'Management'),
('Bob', 'Architect', 15000.00, 'IT'),
('Charlie', 'Manager', 12000.00, 'Marketing'),
('Rahim', 'Uddin', 4500.00, 'Operations');

-- 4. BULK INSERT (Randomized Data)
-- This generates 1000 fake employees
INSERT INTO employees (first_name, last_name, salary, department)
SELECT 
    -- Pick a random first name from the array (Mixed International & Bangladeshi)
    (ARRAY[
        'John','Jane','Mike','Sara','Tom','Emily','Chris','Anna','Steve','Laura',
        'Rahim','Karim','Fatima','Ayesha','Rafiq','Jabbar','Nasrin','Monir','Sultana','Kamal'
    ])[floor(random() * 20 + 1)],
    
    -- Pick a random last name (Mixed International & Bangladeshi)
    (ARRAY[
        'Smith','Johnson','Williams','Jones','Brown','Davis','Miller','Wilson','Moore','Taylor',
        'Ahmed','Hossain','Khan','Chowdhury','Rahman','Islam','Sarker','Uddin','Ali','Hasan'
    ])[floor(random() * 20 + 1)],
    
    -- Generate random salary between 3,000 and 18,000
    ROUND((random() * 15000 + 3000)::numeric, 2),
    
    -- Pick a random department
    (ARRAY['IT', 'HR', 'Management', 'Marketing', 'Sales', 'Operations'])[floor(random() * 6 + 1)]
FROM generate_series(1, 1000);

-- 5. VERIFY
-- Should return 1004 rows
SELECT count(*) as total_employees FROM employees;
\end{lstlisting}

% =======================================================
% SECTION: INTRODUCTION TO PL/pgSQL
% =======================================================
\newpage
\section{Introduction to PL/pgSQL}

PL/pgSQL (Procedural Language/PostgreSQL Structured Query Language) extends SQL by adding procedural constructs. While SQL is declarative (stating \textit{what} you want), PL/pgSQL allows you to state \textit{how} to achieve it using loops, variables, and logic.

\subsection{Anonymous Blocks}
An \textbf{Anonymous Block} is a piece of code that is executed immediately and not stored in the database. In PostgreSQL, this is done using the \texttt{DO} statement.

\textbf{Structure:}
The code body must be enclosed in Dollar Quotes (\texttt{\$\$}).


\begin{lstlisting}[language=SQL]
DO $$ 
DECLARE
    -- Variable declarations
BEGIN
    -- Execution logic
EXCEPTION
    -- Error handling
END $$;
\end{lstlisting}

\subsection{Outputting Text}
PostgreSQL uses \texttt{RAISE NOTICE} to print messages to the console (similar to \texttt{System.out.println} in Java or \texttt{printf} in C).

% =======================================================
% SECTION: VARIABLES \& ARITHMETIC
% =======================================================
\section{Variables and Arithmetic}

Variables must be declared in the \texttt{DECLARE} section. You can initialize them using the \texttt{:=} operator.

\begin{lstlisting}[language=SQL, caption={Geometry Calculation}]
DO $$ 
DECLARE
    v_radius    NUMERIC := 5.5;
    v_pi        CONSTANT NUMERIC := 3.14159;
    v_area      NUMERIC;
BEGIN
    -- Formula: Area = pi * r^2
    v_area := v_pi * (v_radius ^ 2);

    RAISE NOTICE 'Radius: %', v_radius;
    RAISE NOTICE 'Calculated Area: %', v_area; 
END $$;
\end{lstlisting}

% =======================================================
% SECTION: INTERACTING WITH DATA
% =======================================================
\section{Single Row Queries (SELECT INTO)}

To use data from a table inside a PL/pgSQL block, you cannot simply write \texttt{SELECT * FROM table}. The result must be stored in a variable using the \texttt{INTO} keyword.

\textbf{Rule:} The query must return exactly \textbf{one} row. If it returns zero rows or multiple rows, it may raise an error or warning depending on configuration.

\begin{lstlisting}[language=SQL, caption={Fetching Employee Data}]
DO $$ 
DECLARE
    v_emp_id        INT := 1;
    v_full_name     TEXT;
    v_salary        NUMERIC;
BEGIN
    -- 1. Fetch data into variables
    SELECT first_name || ' ' || last_name, salary 
    INTO v_full_name, v_salary
    FROM employees 
    WHERE employee_id = v_emp_id;

    RAISE NOTICE 'Employee: % earns %', v_full_name, v_salary;

    -- 2. Use variables in DML (Update)
    UPDATE employees 
    SET salary = salary + 100 
    WHERE employee_id = v_emp_id;
    
    RAISE NOTICE 'Salary updated.';
END $$;
\end{lstlisting}

% =======================================================
% SECTION: CONTROL STRUCTURES
% =======================================================
\newpage
\section{Control Structures}

\subsection{Conditional Logic (IF/ELSE)}
PL/pgSQL supports \texttt{IF-THEN-ELSIF-ELSE} logic.

\begin{lstlisting}[language=SQL, caption={Salary Evaluation}]
DO $$ 
DECLARE
    v_salary NUMERIC;
BEGIN
    -- Get Alice's salary
    SELECT salary INTO v_salary FROM employees WHERE first_name = 'Alice';

    IF v_salary > 10000 THEN
        RAISE NOTICE 'Status: Executive Level';
    ELSIF v_salary BETWEEN 5000 AND 10000 THEN
        RAISE NOTICE 'Status: Manager Level';
    ELSE
        RAISE NOTICE 'Status: Entry Level';
    END IF;
END $$;
\end{lstlisting}

\subsection{CASE Statement}
The \texttt{CASE} statement is useful for selecting a value based on multiple patterns.

\begin{lstlisting}[language=SQL]
DO $$ 
DECLARE
    v_grade CHAR(1) := 'B';
    v_result TEXT;
BEGIN
    CASE v_grade
        WHEN 'A' THEN v_result := 'Excellent';
        WHEN 'B' THEN v_result := 'Good';
        ELSE v_result := 'Average';
    END CASE;
    
    RAISE NOTICE 'Grade B means: %', v_result;
END $$;
\end{lstlisting}

% =======================================================
% SECTION: LOOPS
% =======================================================
\section{Loops}

\subsection{WHILE Loop}
Executes as long as the condition is true.

\begin{lstlisting}[language=SQL]
DO $$ 
DECLARE 
    v_counter INT := 1;
BEGIN
    WHILE v_counter <= 3 LOOP
        RAISE NOTICE 'Count: %', v_counter;
        v_counter := v_counter + 1;
    END LOOP;
END $$;
\end{lstlisting}

\subsection{FOR Loop (Iterating Query Results)}
This is the most powerful loop in PL/pgSQL. It allows you to iterate over every row returned by a query without explicitly opening or closing cursors.

\begin{lstlisting}[language=SQL, caption={Looping through rows}]
DO $$ 
DECLARE
    r RECORD; -- A generic holder for row data
BEGIN
    RAISE NOTICE '--- IT Department Staff ---';
    
    -- "r" becomes a row for every iteration
    FOR r IN 
        SELECT first_name, salary FROM employees WHERE department = 'IT' 
    LOOP
        RAISE NOTICE 'Name: %, Salary: %', r.first_name, r.salary;
    END LOOP;
END $$;
\end{lstlisting}

% =======================================================
% SECTION: ADVANCED CONCEPTS
% =======================================================
\newpage
\section{Using WITH (CTEs) in Blocks}

You cannot start a \texttt{DO} block with \texttt{WITH}. However, you can use Common Table Expressions (CTEs) inside the SQL statements contained within the block.

\begin{lstlisting}[language=SQL, caption={CTE inside Anonymous Block}]
DO $$ 
DECLARE
    r RECORD;
BEGIN
    -- We want to find employees who earn more than the average
    FOR r IN 
        WITH AvgSal AS (
            SELECT AVG(salary) as val FROM employees
        )
        SELECT e.first_name, e.salary, a.val as average
        FROM employees e, AvgSal a
        WHERE e.salary > a.val
    LOOP
        RAISE NOTICE '% earns % (Above Avg: %)', 
            r.first_name, r.salary, ROUND(r.average, 2);
    END LOOP;
END $$;
\end{lstlisting}

% =======================================================
% SECTION: LAB TASKS
% =======================================================
\newpage
\section{Lab Tasks}

\textbf{Scenario:} You are the database administrator for a startup. You need to perform complex updates and generate reports using PL/pgSQL.


\subsection*{Task 1}
\textbf{Problem:} Find the Min and Max salary without using aggregate functions (\texttt{MIN}/\texttt{MAX}).
\begin{itemize}
    \item Initialize \texttt{v\_min} to a very high number and \texttt{v\_max} to 0.
    \item Loop through all employees.
    \item Compare current salary with variables and update if necessary.
    \item Print results.
\end{itemize}

\ifshowsolutions
\textbf{Solution:}
\begin{lstlisting}[language=SQL]
DO $$ 
DECLARE
    r RECORD;
    v_min NUMERIC := 999999;
    v_max NUMERIC := 0;
BEGIN
    FOR r IN SELECT salary FROM employees LOOP
        
        IF r.salary > v_max THEN
            v_max := r.salary;
        END IF;

        IF r.salary < v_min THEN
            v_min := r.salary;
        END IF;

    END LOOP;
    
    RAISE NOTICE 'Found Max: %, Found Min: %', v_max, v_min;
END $$;
\end{lstlisting}
\fi

\subsection*{Task 2}
Safely remove low-budget departments.
\begin{itemize}
    \item Check if the \texttt{'Interns'} department exists.
    \item If it exists, count the members.
    \item If count $> 0$, delete them and print "Deleted X interns".
    \item If count is 0, print "No interns found to delete".
\end{itemize}

\ifshowsolutions
\textbf{Solution:}
\begin{lstlisting}[language=SQL]
DO $$ 
DECLARE
    v_count INT;
BEGIN
    SELECT COUNT(*) INTO v_count FROM employees WHERE department = 'Interns';

    IF v_count > 0 THEN
        DELETE FROM employees WHERE department = 'Interns';
        RAISE NOTICE 'Deleted % interns', v_count;
    ELSE
        RAISE NOTICE 'No interns found to delete';
    END IF;
END $$;
\end{lstlisting}
\fi

\subsection*{Task 3}
Categorize employees into bands.
\begin{itemize}
    \item Band A: $> 15,000$.
    \item Band B: $10,000 - 15,000$.
    \item Band C: $< 10,000$.
    \item Count how many are in each band using variables inside a loop.
    \item Print the final counts.
\end{itemize}

\ifshowsolutions
\textbf{Solution:}
\begin{lstlisting}[language=SQL]
DO $$ 
DECLARE
    r RECORD;
    v_band_a INT := 0;
    v_band_b INT := 0;
    v_band_c INT := 0;
BEGIN
    FOR r IN SELECT salary FROM employees LOOP
        IF r.salary > 15000 THEN
            v_band_a := v_band_a + 1;
        ELSIF r.salary >= 10000 THEN
            v_band_b := v_band_b + 1;
        ELSE
            v_band_c := v_band_c + 1;
        END IF;
    END LOOP;

    RAISE NOTICE 'Band A: %, Band B: %, Band C: %', v_band_a, v_band_b, v_band_c;
END $$;
\end{lstlisting}
\fi
\subsection*{Task 4}
\textbf{Problem:} Find the next prime number.
\begin{itemize}
    \item Define an input variable \texttt{v\_n}.
    \item Write logic to find the smallest prime number that is \textbf{strictly greater} than \texttt{v\_n}.
    \item Print: "Next prime after [n] is [result]".
\end{itemize}

\ifshowsolutions
\textbf{Solution:}
\begin{lstlisting}[language=SQL]
DO $$ 
DECLARE
    -- ===========================
    -- INPUT SYNTAX
    -- Change this value to test
    -- ===========================
    v_n INT := 14; 
    
    -- Operational variables
    v_candidate INT;
    v_i INT;
    v_is_prime BOOLEAN;
BEGIN
    -- Start checking from the next number
    v_candidate := v_n + 1;

    LOOP
        -- Assume current candidate is prime until proven otherwise
        v_is_prime := TRUE;

        -- 1 and 0 are not primes (though we are checking > n)
        IF v_candidate < 2 THEN
            v_is_prime := FALSE;
        ELSE
            -- Check divisibility from 2 up to square root of candidate
            FOR v_i IN 2..FLOOR(SQRT(v_candidate))::INT LOOP
                IF v_candidate % v_i = 0 THEN
                    v_is_prime := FALSE;
                    EXIT; -- Found a factor, not prime
                END IF;
            END LOOP;
        END IF;

        -- If logic confirms it is prime, print and break
        IF v_is_prime THEN
            RAISE NOTICE 'Next prime after % is %', v_n, v_candidate;
            EXIT;
        END IF;

        -- Otherwise, increment and try the next number
        v_candidate := v_candidate + 1;
    END LOOP;
END $$;
\end{lstlisting}
\fi

\end{document}